\chapter{Resumen Ejecutivo}

El proyecto \textbf{iComanda} consiste en el desarrollo de una aplicación móvil
y web orientada a restaurantes pequeños y medianos que necesitan digitalizar la
gestión de comandas, mesas, menú, cocina, pagos, gastos, inventario y reportes
operativos. La solución reemplaza el registro manual en papel por un flujo
centralizado donde meseros, administradores y cocina trabajan sobre la misma
información en tiempo real.

El sistema fue construido con \textbf{React Native, Expo Router, TypeScript y
Supabase}. La aplicación móvil permite operar desde dispositivos Android, mientras
que las rutas públicas web permiten que los clientes consulten la carta digital
desde un navegador mediante códigos QR, sin instalar una aplicación ni crear una
cuenta tradicional.

\section{Problema Abordado}

En muchos restaurantes pequeños y medianos, los pedidos todavía se registran en
papel o mediante mensajes informales entre meseros y cocina. Esto produce errores
en los totales, pérdida de comandas, dificultad para conocer el estado real de una
mesa y poca visibilidad financiera al final del turno. Además, la administración
no cuenta con información inmediata sobre ventas, gastos, platos más vendidos o
disponibilidad de insumos.

iComanda aborda esta necesidad mediante una plataforma que integra operación diaria
y control administrativo. El mesero registra pedidos desde el celular, cocina ve
las órdenes pendientes y las marca como listas, el administrador gestiona menú,
usuarios, mesas, pagos QR, gastos, inventario y reportes, y el cliente puede ver
la carta o llamar al mesero desde la mesa.

\section{Solución Implementada}

La solución implementada incluye los siguientes módulos principales:

\begin{itemize}[leftmargin=1.5em]
  \item \textbf{Autenticación y roles:} acceso diferenciado para \texttt{admin},
        \texttt{mesero} y \texttt{chef} mediante Supabase Auth y perfiles en la
        tabla \texttt{profiles}.
  \item \textbf{Gestión de mesas:} administración de mesas, capacidad, estado
        activo y generación de enlaces QR por mesa.
  \item \textbf{Gestión de menú:} creación, edición, disponibilidad, imágenes y
        marcado de platos agotados.
  \item \textbf{Toma de pedidos:} registro de órdenes con mesa, ítems, cantidades,
        método de pago y cálculo automático del total.
  \item \textbf{Cocina:} panel para que el rol \texttt{chef} visualice pedidos en
        preparación y marque órdenes como \texttt{lista}.
  \item \textbf{Tiempo real:} actualización de pedidos, mesas, llamadas del cliente
        y disponibilidad mediante Supabase Realtime.
  \item \textbf{Carta digital pública:} menú accesible desde navegador por ruta
        \texttt{/menu/[slug]}.
  \item \textbf{Vista pública de cliente en mesa:} ruta \texttt{/table/[slug]/[tableId]}
        para ver el menú y llamar al mesero desde el QR de mesa, sin crear cuenta.
  \item \textbf{Pagos:} soporte para efectivo, QR y tarjeta, con confirmación de
        pagos QR y registro de \texttt{paid\_at} para reportes.
  \item \textbf{Gastos e inventario:} registro de egresos, materiales, reposiciones,
        consumos y control básico de stock.
  \item \textbf{Reportes:} métricas financieras, ranking de platos vendidos y ventas
        por mesero.
\end{itemize}

\section{Arquitectura Técnica}

El proyecto utiliza una arquitectura \textbf{Backend as a Service (BaaS)}. La
aplicación cliente consume directamente los servicios de Supabase, evitando la
necesidad de implementar un backend intermedio propio.

\begin{itemize}[leftmargin=1.5em]
  \item \textbf{Frontend móvil y web:} React Native, Expo SDK 54, Expo Router 6.
  \item \textbf{Lenguaje:} TypeScript.
  \item \textbf{Backend administrado:} Supabase con PostgreSQL, Auth, Storage,
        Realtime, funciones SQL y políticas RLS.
  \item \textbf{Diseño de interfaz:} componentes propios inspirados en Material
        Design 3.
  \item \textbf{Control de versiones:} Git y GitHub, usando ramas auxiliares de
        desarrollo que se integraron a la rama principal y luego se eliminaron.
\end{itemize}

\section{Metodología y Resultado General}

El desarrollo se organizó bajo una metodología ágil basada en Scrum, dividido en
\textbf{3 sprints de 2 semanas}. Cada sprint entregó incrementos funcionales:
fundación técnica y autenticación, operación de pedidos y administración, y
funcionalidades finales de tiempo real, carta digital, cocina, inventario y
reportes.

Como resultado, se obtuvo una solución funcional que centraliza la operación de un
restaurante en una sola plataforma, mejora la trazabilidad de pedidos y pagos,
reduce errores del registro manual y proporciona información útil para la toma de
decisiones administrativas.
